\chapter{Proposed Methodology}
\label{chap:methodology}

\section{Introduction}

This chapter describes how the comparison announced in Chapter~\ref{chap:introduction} is
carried out. The question is narrow on purpose: on one fixed corpus, one fixed encoder and one
fixed training protocol, what changes when only the prediction head and the loss function
change? Answering it requires that everything else stay identical, so most of this chapter is
concerned with what was held fixed and with how that was verified.

Section~\ref{sec:meth-overview} states the design and the two experimental tracks.
Section~\ref{sec:meth-data} describes the BAREC corpus, the official splits and the checks run
on the loaded data. Section~\ref{sec:meth-input} describes the D3Tok input representation.
Section~\ref{sec:meth-heads} describes the encoder and the three prediction heads.
Section~\ref{sec:meth-losses} gives the exact form of every loss function, together with the
numerical safeguards and the tests that verify each implementation.
Section~\ref{sec:meth-coral-init} presents the cut-point initialisation for CORAL, which is a
contribution of this thesis. Section~\ref{sec:meth-protocol} states the training protocol and
every point at which the implementation departs from the published BAREC script.
Section~\ref{sec:meth-methods} lists the methods compared and the ablations.
Section~\ref{sec:meth-metrics} defines the evaluation metrics.
Section~\ref{sec:meth-frozen} describes the frozen-backbone track that answers the
CPU-friendly requirement of the research brief, and Section~\ref{sec:meth-impl} reports the
implementation, the hardware and the reproducibility record.

\section{Experimental Design}
\label{sec:meth-overview}

The experiments follow the pipeline of the BAREC benchmark \cite{elmadani2025barec}, whose
training script is treated here as the specification of the protocol. The corpus, the input
representation, the encoder, the optimiser, the learning rate, the number of epochs, the batch
size, the schedule, the random seed and the checkpoint selection rule are the same for every
method. Two components vary: the prediction head, which turns a sentence representation into
either 19 scores, one continuous value or 18 cumulative scores, and the loss function applied
to that output. This isolation is the reason the comparison can be attributed to the ordinal
mechanism rather than to a difference in tuning.

Two tracks are reported.

\begin{enumerate}
\item \textbf{Fine-tuning track (main results).} The whole encoder is fine-tuned, exactly as in
the BAREC benchmark. This track is the only one whose numbers are comparable with the published
tables, and it carries the answer to the research question.

\item \textbf{Frozen-backbone track.} The encoder is run once in inference mode, its sentence
features are cached, and only a linear head is trained on those cached features. Head training
takes minutes on a laptop processor, which satisfies the CPU-friendly setting requested in the
research brief, and it is cheap enough to repeat over five random seeds, so every number in
that track is reported with a standard deviation. Because the encoder never adapts to the task,
these numbers are far below the fine-tuning track and are never compared with it or with the
published tables.
\end{enumerate}

A deliberate limitation of the fine-tuning track should be stated at the outset: each method is
trained once, with the seed fixed by the reference protocol at 42. The frozen track is the part
of this work that reports variance.

\section{Data}
\label{sec:meth-data}

\subsection{The BAREC Corpus}

All experiments use the Balanced Arabic Readability Evaluation Corpus
\cite{elmadani2025barec}, obtained from the public repository \texttt{\datasetName}. The corpus
contains 69{,}441 manually annotated Arabic sentences over 19 ordered readability levels, named
after the Abjad ordering of the Arabic letters (1-\textit{alif} to 19-\textit{qaf}) and ranging
from kindergarten to postgraduate reading difficulty. Sentences carry the fine-grained label in
the column \texttt{Readability\_Level\_19} and the three collapsed labels
\texttt{Readability\_Level\_7}, \texttt{Readability\_Level\_5} and
\texttt{Readability\_Level\_3}.

The corpus is split at the document level, so that no document contributes sentences to more
than one split. Table~\ref{tab:meth-splits} reports the sizes.

\begin{table}[H]
\centering
\caption{BAREC splits as used in this thesis. The sizes are asserted against Table~3 of
Elmadani et al.~\cite{elmadani2025barec} before any training starts.}
\label{tab:meth-splits}
\begin{tabular}{lrrr}
\toprule
\textbf{Split} & \textbf{Documents} & \textbf{Sentences} & \textbf{Words} \\
\midrule
Train       & 1{,}518 & 54{,}845 & 832{,}743 \\
Development &   194 &  7{,}310 & 101{,}364 \\
Test        &   210 &  7{,}286 & 105{,}264 \\
\midrule
All         & 1{,}922 & 69{,}441 & 1{,}039{,}371 \\
\bottomrule
\end{tabular}
\end{table}

Labels are shifted to $y \in \{0,\dots,18\}$ during training and shifted back only when results
are reported, which follows the reference script and keeps the index arithmetic in the losses
simple.

\subsection{Checks Run on the Loaded Data}

Four assertions run before any model is built, and the program stops rather than continue if
one of them fails.

\begin{enumerate}
\item The three split sizes equal the published values in Table~\ref{tab:meth-splits}.
\item The D3Tok column exists, and the clitic marker \texttt{+} occurs in more than half of a
2{,}000-sentence sample of the training split. A silent preprocessing failure would leave plain
text in that column, and this check catches it.
\item The 7-, 5- and 3-level collapse tables of the reference script are applied to the
19-level gold labels and compared, sentence by sentence, with the corpus's own collapsed
columns on all three splits. They agree everywhere, so the collapsed accuracies reported in
Chapter~\ref{chap:results} cannot be the product of a misremembered mapping.
\item The development and test splits each contain at least one sentence of every one of the 19
levels, which the quadratic weighted kappa computation requires to be meaningful.
\end{enumerate}

One correction to the published script is needed: the released dataset names its middle split
\texttt{dev}, while the script reads a split named \texttt{validation}. The implementation
accepts either name and then verifies the size, so the correct split is used.

\section{Input Representation}
\label{sec:meth-input}

\subsection{D3Tok}

The corpus ships several input variants of every sentence. This thesis uses D3Tok, the
morphologically segmented variant produced by the corpus authors with CAMeL Tools
\cite{obeid2020camel}: each word is separated into its base and its clitics, and every clitic
boundary is marked with a \texttt{+} character. D3Tok is read as a column of the corpus, so the
segmentation is the corpus authors' own and is identical for every method, which removes the
preprocessing tool as a source of variation between methods.

Arabic attaches function words to their host word in writing. A single written token may fuse a
conjunction, the definite article and a noun; in Buckwalter transliteration such a token
appears as \texttt{wAlfDp} and D3Tok renders it \texttt{w+ Al+ fDp}, exposing the conjunction
\textit{wa+}, the article \textit{al+} and the noun. The BAREC annotation guidelines list the
presence and the type of clitics among the markers of readability level
\cite{habash2025guidelines}, so exposing them to the encoder is linguistically motivated and
not only a tokenisation convenience.

Two empirical reasons support the choice. First, AraBERTv2 was pretrained on morphologically
segmented Arabic \cite{antoun2020arabert}, so segmented input matches the distribution seen
during pretraining. Second, in the benchmark itself, the pairing of AraBERTv2 with D3Tok input
raises development QWK from 76.3 to 80.0 relative to plain word input, and that pairing is the
configuration the benchmark carries into its loss function comparison
\cite{elmadani2025barec}. Using it here is what makes the reproduction check in
Chapter~\ref{chap:results} meaningful.

\subsection{Tokenisation}

Sentences are tokenised with the AraBERTv2 tokenizer, truncated at 512 subword tokens, and
padded dynamically to the longest sequence in each batch. Truncation length and dynamic padding
both follow the reference script. Dynamic padding changes only how much padding the encoder
processes, never the content of a sentence.

\section{Encoder and Prediction Heads}
\label{sec:meth-heads}

\subsection{Encoder}

The encoder is AraBERTv2 (\texttt{aubmindlab/bert-base-arabertv2})
\cite{antoun2020arabert}: twelve transformer layers, twelve attention heads and a hidden size
of 768, pretrained on approximately 8.6 billion words of Arabic. The same checkpoint, at the
same repository revision, is used by every method in the fine-tuning track, and all of its
parameters are updated during training.

The sentence representation is the encoder's pooled output, that is, the first-token
representation passed through the pretrained pooling layer with its hyperbolic tangent
activation. A dropout layer with the rate stored in the model configuration, $0.1$, is applied
to it. This is the representation that the standard sequence classification wrapper uses, so
the classification and regression heads take it unchanged; the CORAL head is built on the same
representation so that the three heads differ in their output structure only.

\subsection{Three Heads}

Let $\mathbf{h} \in \mathbb{R}^{768}$ be the pooled, dropped-out representation and let $K=19$.

\paragraph{Classification head.} A linear layer produces one score per level,
\begin{equation}
\mathbf{z} = W\mathbf{h} + \mathbf{b}, \qquad W \in \mathbb{R}^{19\times768},\quad
\mathbf{b}\in\mathbb{R}^{19},
\end{equation}
which a softmax turns into a distribution $\mathbf{p}$ over the 19 levels. Prediction is
$\hat{y} = \arg\max_k p_k$. Cross-entropy, OLL, EMD and the four soft-label losses all use this
head; they differ only in the loss applied to $\mathbf{p}$.

\paragraph{Regression head.} A linear layer produces one unbounded scalar,
\begin{equation}
\hat{s} = \mathbf{w}^{\top}\mathbf{h} + b, \qquad \mathbf{w}\in\mathbb{R}^{768},\quad
b\in\mathbb{R}.
\end{equation}
Prediction rounds to the nearest level and clips to the valid range,
$\hat{y} = \mathrm{clip}(\mathrm{rint}(\hat{s}), 0, 18)$, where $\mathrm{rint}$ rounds halves
to the nearest even integer. The mean squared error method and the Soft-QWK method use this
head.

\paragraph{CORAL head.} The CORAL head \cite{cao2020coral} replaces the 19-way decision by
$K-1=18$ ordered binary questions, ``is the level greater than $k$?''. All 18 questions share a
single projection vector and differ only in a free bias, called a cut-point:
\begin{equation}
z_k = \mathbf{w}^{\top}\mathbf{h} + b_k, \qquad k = 0,\dots,17 ,
\label{eq:coral-logits}
\end{equation}
implemented as a bias-free linear layer of output size one, giving a single readability score
$s = \mathbf{w}^{\top}\mathbf{h}$, plus a learnable vector of 18 cut-points added to it. After
a sigmoid, $\sigma(z_k)$ estimates $P(y>k)$. Prediction counts the questions answered yes,
\begin{equation}
\hat{y} = \sum_{k=0}^{17} \mathbb{1}\left[\sigma(z_k) > 0.5\right].
\label{eq:coral-decode}
\end{equation}

Because the projection is shared, the 18 estimated probabilities are ordered whenever the
cut-points are ordered, so the model cannot claim that a sentence exceeds level 9 while denying
that it exceeds level 5. Cao et al.\ prove that the minimiser of the CORAL loss has ordered
cut-points \cite{cao2020coral}; that is a statement about the optimum, not a guarantee about
every intermediate model, so whether the cut-points are in fact ordered after each epoch is
logged during training and reported in Chapter~\ref{chap:results}.

The three heads add 14{,}611, 769 and 786 parameters respectively on top of the encoder, so the
capacity added by the head is negligible against the roughly 135 million parameters of AraBERTv2 and
cannot explain any difference between methods.

Table~\ref{tab:meth-heads} summarises the heads and their decoding rules.

\begin{table}[H]
\centering
\footnotesize
\caption{Prediction heads, their outputs and their decoding rules.}
\label{tab:meth-heads}
\begin{tabular}{p{2.6cm} p{1.9cm} p{4.3cm} p{4.0cm}}
\toprule
\textbf{Head} & \textbf{Output} & \textbf{Used by} & \textbf{Decoding} \\
\midrule
Classification & $\mathbb{R}^{19}$ & CE, OLL, EMD, SOFT ($\beta=2,3,4,5$) &
$\arg\max_k p_k$ \\
Regression & $\mathbb{R}$ & Regression (MSE), MSE $\rightarrow$ Soft-QWK &
round to nearest level, clip to $[0,18]$ \\
CORAL & $\mathbb{R}^{18}$ & CORAL and its ablation &
count of $\sigma(z_k) > 0.5$ \\
\bottomrule
\end{tabular}
\end{table}

\section{Loss Functions}
\label{sec:meth-losses}

Throughout this section $y\in\{0,\dots,18\}$ is the gold level of one sentence, $\mathbf{p}$ is
the predicted distribution from the classification head, $\hat{s}$ is the output of the
regression head, and $N$ is the number of sentences in the batch. Every loss is a mean over the
batch, which is the reduction used by the reference script.

\subsection{Cross-Entropy}

\begin{equation}
\mathcal{L}_{\text{CE}} = -\log p_y .
\end{equation}

Cross-entropy is the ordinal-blind baseline. Its value depends only on the probability given to
the correct level, so moving the remaining mass from an adjacent level to the most distant
level does not change it. Every other loss in this comparison exists to remove that property.

\subsection{Regression with Mean Squared Error}

\begin{equation}
\mathcal{L}_{\text{MSE}} = (\hat{s}-y)^2 .
\end{equation}

The penalty grows with the square of the error, so a prediction ten levels away costs one
hundred times a prediction one level away. Ordinal information enters through the arithmetic of
the label rather than through an explicit distance matrix. This is the method that obtained the
best QWK in the benchmark \cite{elmadani2025barec}.

\subsection{Ordinal Log Loss}

Ordinal log loss \cite{castagnos2022ordinal} keeps the classification head and weights the
probability given to each wrong level by its distance from the gold level:
\begin{equation}
\mathcal{L}_{\text{OLL}} = -\sum_{k=0}^{18} |y-k|^{\alpha}\,\log\left(1-p_k\right),
\qquad \alpha = 1.5 .
\end{equation}
The gold level contributes nothing, since its weight $|y-y|^{\alpha}$ is zero; the loss is
driven entirely by where the incorrect mass sits. The exponent $\alpha$ controls how sharply
distant mass is punished relative to adjacent mass, and $\alpha=1.5$ is the value that performed
best among $\{0.5, 1, 1.5, 2\}$ in the benchmark, which is why it is the value used here.

\subsection{Earth Mover's Distance}

The EMD loss \cite{hou2017emd} compares cumulative distributions rather than probabilities.
With $P_k=\sum_{j\leq k}p_j$ and the gold cumulative distribution
$T_k = \mathbb{1}[y \leq k]$,
\begin{equation}
\mathcal{L}_{\text{EMD}} = \sum_{k=0}^{18}\left(P_k - T_k\right)^2 .
\end{equation}
The sum is not divided by the number of levels, again following the reference implementation.
Probability mass that sits near the gold level produces a short mismatch between the two
cumulative curves, and mass that sits far from it produces a long one.

\subsection{Soft Labels}

The soft-label loss \cite{bertinetto2020soft} replaces the one-hot target by a distribution
that decays exponentially with distance from the gold level, and then applies cross-entropy
against that target:
\begin{equation}
q_i = \frac{e^{-\beta|y-i|}}{\sum_{n=0}^{18} e^{-\beta |n-i|}}, \qquad
\mathcal{L}_{\text{SOFT}} = -\sum_{i=0}^{18} q_i \log p_i .
\end{equation}
The normalisation runs over the levels $n$ for a fixed $i$, as in the reference
implementation, so the target is not exactly a probability vector over $i$; it is a
distance-weighted target whose scale depends on $\beta$. The smoothing parameter is evaluated
at $\beta \in \{2,3,4,5\}$, the grid used in the benchmark. A large $\beta$ concentrates the
target on the gold level and approaches cross-entropy; a small $\beta$ spreads credit onto
neighbouring levels.

\subsection{CORAL}

For the CORAL head, the target of question $k$ is $t_k = \mathbb{1}[y>k]$ and the loss is the
binary cross-entropy of the 18 questions, summed over questions and averaged over the batch:
\begin{equation}
\mathcal{L}_{\text{CORAL}} = \frac{1}{N}\sum_{n=1}^{N}\sum_{k=0}^{17}
\Big[-t_k^{(n)}\log\sigma\big(z_k^{(n)}\big)
-\big(1-t_k^{(n)}\big)\log\big(1-\sigma\big(z_k^{(n)}\big)\big)\Big].
\label{eq:coral-loss}
\end{equation}
Summing rather than averaging over the 18 questions is what Cao et al.\ specify
\cite{cao2020coral}. The choice matters: averaging would divide every gradient reaching the
shared projection and the cut-points by 18, which at a fixed learning rate would shrink the
distance those parameters can travel by the same factor. Averaging over the batch is the
reduction used by all the other losses here, so the two conventions are combined as written in
Equation~\ref{eq:coral-loss}.

\subsection{Soft-QWK, and Why It Needs Two Phases}

The metric used to judge every method, quadratic weighted kappa, cannot be optimised directly,
because it is computed from counts of discrete predictions and counting has no useful gradient.
The STBW system \cite{trigui2025stbw} replaces each discrete step by a smooth one, following
earlier work on differentiable kappa objectives \cite{delatorre2018wkl} and on turning a scalar
prediction into a soft distribution \cite{diaz2019softlabels}. This thesis implements the same
objective.

First, the scalar prediction is clamped into the valid range and spread over the 19 levels with
a Gaussian window of width $\sigma=1$, normalised to sum to one:
\begin{equation}
\hat{y} = \mathrm{clip}(\hat{s},0,18), \qquad
P_i \propto \exp\left(-\frac{(i-\hat{y})^2}{2\sigma^2}\right) .
\end{equation}
The window is centred on the continuous prediction rather than on the rounded one, because a
rounded centre would again have zero gradient. Second, with $\mathbf{T}$ the matrix of one-hot
gold labels and $\mathbf{P}$ the matrix of soft predictions for the batch, a soft agreement
matrix and a chance agreement matrix are accumulated,
\begin{equation}
O = \mathbf{T}^{\top}\mathbf{P}, \qquad
E = \frac{\mathbf{o}_{\text{gold}}\,\mathbf{o}_{\text{pred}}^{\top}}{N},
\end{equation}
where $\mathbf{o}_{\text{gold}}$ and $\mathbf{o}_{\text{pred}}$ are the row and column totals
of $O$. Third, with the quadratic penalty $W_{ij}=(i-j)^2$, the differentiable kappa and the
loss are
\begin{equation}
\kappa_{\text{soft}} = 1 - \frac{\langle W,O\rangle}{\langle W,E\rangle + \varepsilon},
\qquad
\mathcal{L}_{\text{Soft-QWK}} = 1-\kappa_{\text{soft}},
\qquad \varepsilon = 10^{-9}.
\end{equation}

This objective has a property no other loss in the comparison has: it is a statistic of the
whole batch, not a sum of per-sentence terms. A model can therefore lower it by placing every
prediction near the mean level of the batch, which keeps the soft agreement matrix plausible
without reading any sentence. The STBW authors report exactly this failure and avoid it by
training in two phases: mean squared error first, then Soft-QWK from that checkpoint
\cite{trigui2025stbw}. The same two-phase schedule is used here. Phase one trains with mean
squared error for three epochs and keeps the checkpoint with the lowest development loss; phase
two continues from that checkpoint with the Soft-QWK objective for three epochs, with a fresh
optimiser and a fresh learning rate schedule. The total budget stays at six epochs, so the
method is not given more training than the others, but its schedule is different and its
results are therefore reported as a separate group throughout Chapter~\ref{chap:results}.

Three values of this method are not specified by the STBW paper and are choices made here: the
Gaussian width $\sigma=1$, the position of the phase boundary after epoch three, and the
numerical guard $\varepsilon$. The phase boundary follows the STBW observation that their
warm-up checkpoint was the third epoch.

\subsection{Numerical Safeguards}

Two losses take a logarithm of a quantity that a softmax can drive to a boundary. In OLL, a
predicted probability of exactly one at some level would make $\log(1-p_k)$ infinite, and at
the gold level that infinity would be multiplied by a zero weight, producing an undefined
value; the probability is therefore clamped just below one before the logarithm. In the
soft-label loss the probability is clamped just above zero before $\log p_i$. Both clamps use
the smallest representable margin of the working precision, so they agree with the literal
formula wherever the literal formula is finite. All losses upcast the model output to
single precision before any softmax, logarithm or sigmoid, so that mixed-precision training
does not change the arithmetic of the loss itself.

\subsection{Verification of the Loss Implementations}

Each loss is checked against an independent implementation before any training run, on the same
random logits in double precision. Cross-entropy, OLL, EMD and the four soft-label losses are
compared against literal element-by-element loops written in the form used by the reference
script; CORAL is compared against a closed-form numerically stable expansion of the binary
cross-entropy; and Soft-QWK, evaluated with a very narrow Gaussian so that its soft
distribution approaches a one-hot vector, is compared against $1-\kappa$ computed by
\texttt{scikit-learn} on hard predictions. Every check agrees to better than $10^{-9}$, and the
program refuses to train if any of them fails. The metric implementations are verified in the
same way, as described in Section~\ref{sec:meth-metrics}.

\section{Cut-Point Initialisation for CORAL}
\label{sec:meth-coral-init}

The 18 cut-points of the CORAL head are ordinary parameters, and the natural default is to
start them at zero. Under the BAREC protocol that default is a problem, and the reason is
arithmetic rather than statistical.

At convergence, cut-point $b_k$ has to sit near the value at which the shared score $s$
separates level $k$ from level $k+1$. Since the 19 levels are very unevenly populated, those
values are spread over a wide interval: the fitted spread between the highest and the lowest
cut-point in this work is \cutSpreadCoral\ (Chapter~\ref{chap:results}). Now consider how far a
parameter can travel under the protocol. Training runs for six epochs of \stepsPerEpoch\
optimiser steps, that is \gpuTotalSteps\ steps, at a learning rate of $5\times10^{-5}$ decayed
linearly to zero. The sum of the learning rate over all steps is \lrSumEncoder. With an
optimiser whose per-step update is bounded in magnitude by the learning rate, as AdamW's
normalised update effectively is \cite{loshchilov2019adamw}, no single parameter can move much
further than \lrSumEncoder\ in total. A cut-point that starts at zero cannot reach a value near
$-6$ or $+5$, no matter how clear the training signal is.

The initialisation used here removes the problem instead of spending the budget on it. The
cumulative label frequencies of the \emph{training} split are converted into the cut-point that
reproduces them before any sentence is seen:
\begin{equation}
b_k = \log \frac{\hat{P}(y>k)}{1-\hat{P}(y>k)}, \qquad
\hat{P}(y>k) = \frac{1}{N_{\text{train}}}\sum_{n=1}^{N_{\text{train}}}\mathbb{1}\left[y^{(n)}>k\right],
\label{eq:coral-init}
\end{equation}
with the frequency clipped to $[10^{-4}, 1-10^{-4}]$ so that the logit stays finite. Since
$\hat{P}(y>k)$ is non-increasing in $k$ by construction, the initial cut-points are ordered,
which the implementation also asserts. At initialisation the head therefore predicts the
training label distribution with a score of zero, and the training budget is spent on learning
the projection that separates sentences rather than on transporting biases across the real
line. Only training labels are used, so the development and test splits are untouched by the
initialisation.

This is a head initialisation, not a change of the pipeline, and it is the one element of the
method that is neither in the BAREC script nor in the two shared-task papers. Two experiments
in Chapter~\ref{chap:results} test the argument rather than assume it: the distance the
cut-points actually travel during training is measured, and an ablation raises the head
learning rate to $10^{-3}$, which multiplies the budget by twenty while leaving the encoder at
the protocol value.

\section{Training Protocol}
\label{sec:meth-protocol}

\subsection{Hyper-parameters}

Table~\ref{tab:meth-hparams} lists the protocol. Every value is the one used by the BAREC
script, including those that the script leaves at the library default: no warm-up, no weight
decay, a linear decay schedule and gradient clipping at a norm of one.

\begin{table}[H]
\centering
\footnotesize
\caption{Training protocol, identical for every method in the fine-tuning track.}
\label{tab:meth-hparams}
\begin{tabular}{p{6.2cm} p{7.2cm}}
\toprule
\textbf{Setting} & \textbf{Value} \\
\midrule
Encoder & \texttt{aubmindlab/bert-base-arabertv2} \\
Input & D3Tok column of the corpus \\
Maximum sequence length & 512 subword tokens, dynamic padding \\
Optimiser & AdamW, $\beta_1=0.9$, $\beta_2=0.999$, $\epsilon=10^{-8}$ \\
Learning rate & $5\times10^{-5}$ \\
Schedule & linear decay to zero, no warm-up \\
Weight decay & 0 \\
Gradient clipping & maximum norm 1.0 \\
Epochs & 6 \\
Effective training batch & 64 sentences (\microBatch\ per step, \gradAccum\ accumulation steps) \\
Evaluation batch & \evalBatch \\
Optimiser steps per epoch & \stepsPerEpoch \\
Mixed precision & fp16 autocast, single-precision master weights \\
Random seed & 42 \\
Checkpoint selection & lowest development loss over the six epochs, no early stopping \\
\bottomrule
\end{tabular}
\end{table}

\subsection{Checkpoint Selection}

The benchmark selects the epoch with the lowest development loss, keeps no other criterion and
uses no early stopping; that rule is used here for every method, including the ones whose
source papers use a different rule. This matters for the comparison, because the development
loss of each method is measured in that method's own loss, so the rule is not a metric shared
between methods but a per-method criterion. The two shared-task systems both select on
development QWK instead \cite{nazzal2025zai,trigui2025stbw}, which raises the question of how
much of their advantage comes from the selection rule. Section~\ref{sec:meth-methods} therefore
declares a selection ablation that evaluates the highest-QWK epoch of every method as well.

Reporting the numbers of the selected epoch requires care. The selected model is saved
explicitly at the moment it is identified, the saved weights are reloaded for evaluation, and
the reloaded development metrics are compared with what was logged during training at that same
epoch. The comparison is recorded in every result file and is reported in
Chapter~\ref{chap:results}; it is the check that the reported numbers belong to the selected
epoch and not to some later one.

\subsection{Differences from the Published Script}

Three implementation differences are declared here because they are visible in the arithmetic,
and one because it is a correction. None of them changes a hyper-parameter, and all of them are
identical across methods.

\begin{enumerate}
\item \textbf{Batch composition.} A batch of 64 sentences at length 512 does not fit in the 16
gigabytes of the available accelerator, so the batch is formed as \microBatch\ sentences per
forward pass with \gradAccum\ accumulation steps. The 54{,}845 training sentences then give
\stepsPerEpoch\ optimiser steps per epoch, the same number as a true batch of 64. Each
optimiser step therefore averages the gradient of 64 sentences, with one exception: the last
step of each epoch covers the remaining 61 sentences as four micro-batches of 16, 16, 16 and
13 sentences, so its per-sentence weighting differs slightly from the weighting a single batch
of 61 sentences would produce.

\item \textbf{Mixed precision.} Training uses fp16 autocast with single-precision master
weights, which the reference script does not; it was adopted to fit the available time budget.
Every loss is computed in single precision, as described above. Mixed precision can cause the
optimiser to skip a step whose gradients overflow, and the number of such skipped steps was not
recorded, so it cannot be ruled out that a small number of updates were dropped.

\item \textbf{Batch composition for Soft-QWK.} Gradient accumulation is equivalent to a larger
batch only for losses that are sums over sentences. Soft-QWK is a statistic of the whole batch,
so accumulation would compute it on 16 sentences at a time instead of 64. Its second phase is
therefore trained with a true batch of 64 sentences and gradient checkpointing, which trades
computation for memory without changing the arithmetic.

\item \textbf{Methods the script does not implement.} The published script contains no
regression trainer, and it routes the name CORAL to an ordinary classification trainer. The
regression method here follows the description in the benchmark paper, and CORAL follows Cao et
al.\ \cite{cao2020coral}. The weighted kappa loss present in the script is not run, on the
supervisor's instruction, since it appears in no table of the benchmark paper and is neither
CORAL nor Soft-QWK.
\end{enumerate}

\section{Methods Compared and Ablations}
\label{sec:meth-methods}

Table~\ref{tab:meth-methods} lists the eleven training runs. Nine form the comparable group, in
which the protocol is identical and only the head and the loss change. Soft-QWK is reported as
its own group because its schedule differs, as explained above. The two ablations answer
specific questions and are never mixed into the comparable group.

\begin{table}[H]
\centering
\footnotesize
\caption{The eleven runs of the fine-tuning track.}
\label{tab:meth-methods}
\begin{tabular}{p{3.4cm} p{1.7cm} p{2.0cm} p{5.6cm}}
\toprule
\textbf{Run} & \textbf{Head} & \textbf{Group} & \textbf{Role} \\
\midrule
CE & classification & comparable & ordinal-blind baseline \\
Regression (MSE) & regression & comparable & best benchmark loss \\
OLL, $\alpha=1.5$ & classification & comparable & distance-weighted loss \\
EMD & classification & comparable & cumulative distribution loss \\
SOFT, $\beta=2,3,4,5$ & classification & comparable & distance-weighted targets \\
CORAL & CORAL & comparable & ordinal head, the ZAI approach \\
MSE $\rightarrow$ Soft-QWK & regression & two-phase & metric-aligned, the STBW approach \\
CORAL, head lr $10^{-3}$ & CORAL & ablation & tests the cut-point travel argument \\
\bottomrule
\end{tabular}
\end{table}

Two further ablations need no additional training, because they re-evaluate saved models.

\paragraph{Decision threshold for CORAL.} The ZAI system tunes a single decision threshold on
the development set instead of using $0.5$ \cite{nazzal2025zai}. The saved CORAL model is
therefore decoded again with every threshold on the grid $0.05, 0.06, \dots, 0.95$, the value
with the highest development QWK is kept, and the test numbers under that threshold are
reported. Its development numbers are optimistic by construction, since the threshold was
chosen on that split.

\paragraph{Selection rule.} As announced above, the highest development QWK epoch of every
method is also saved and evaluated. Here too the development numbers are optimistic by
construction and only the test column can be read as a result.

\section{Evaluation Metrics}
\label{sec:meth-metrics}

Six metrics come from the benchmark, and three more are added.

\begin{itemize}
\item \textbf{QWK}, quadratic weighted kappa \cite{cohen1968kappa}, computed with
\texttt{scikit-learn} with quadratic weights. It is the primary metric of the benchmark and of
the shared task, and it is the metric on which all comparisons in this thesis are based. It
rewards agreement while penalising an error by the square of its size, which is exactly the
ordinal property the losses try to exploit.
\item \textbf{Acc19}, the proportion of exactly correct levels.
\item \textbf{$\pm$1 Acc19}, the proportion of predictions within one level of the gold level.
\item \textbf{Dist}, the mean absolute distance between predicted and gold level, in levels.
\item \textbf{Acc7, Acc5, Acc3}, exact accuracy after collapsing predictions and gold labels to
the coarser 7-, 5- and 3-level schemes with the corpus's own mappings, verified as described in
Section~\ref{sec:meth-data}.
\item \textbf{far$\geq$3} and \textbf{far$\geq$5}, the proportion of predictions that are at
least three and at least five levels away from the gold level. These two are \emph{not} BAREC
metrics; they are introduced in this thesis to answer the part of the research brief that asks
for the error distance, and specifically to separate a near miss from a far miss. $\pm$1
accuracy counts the near misses and these two count the tails. They are labelled as this
thesis's own wherever they appear.
\item \textbf{Macro-F1} over the 19 levels, reported as a diagnostic of whether a method uses
the whole label range or concentrates on the frequent levels.
\end{itemize}

The metric code is tested before use: a perfect predictor must give QWK, Acc19 and all
collapsed accuracies of one and a distance of zero; a predictor shifted by exactly one level
must give an Acc19 of zero, a $\pm$1 accuracy of one and a distance of exactly one; and
distance, $\pm$1 accuracy and far$\geq$3 are recomputed directly from the absolute differences
and compared with the library results. The three collapse tables are also checked to be onto
the 7, 5 and 3 target levels.

\section{Frozen-Backbone Track}
\label{sec:meth-frozen}

The research brief asks for CPU-friendly training settings. Fine-tuning 135 million parameters
does not meet that description on any processor, so the requirement is met by a second track in
which the encoder is not trained at all.

The encoder is run once over the corpus in inference mode and two sentence representations are
cached for every sentence: the pooled output, as in the fine-tuning track, and the mean of the
final-layer token representations over the non-padding tokens. Only a linear head is then
trained on the cached features, which takes a few minutes per run on a laptop processor; the
\cpuSeeds\ runs of every method together take \cpuTotalMinutes\ minutes.

The head hyper-parameters cannot be the protocol's. A linear head on frozen features, trained
for six epochs at $5\times10^{-5}$, does not converge: the travel budget argument of
Section~\ref{sec:meth-coral-init} applies to every head parameter, not only to CORAL's
cut-points. The head is therefore trained for \cpuEpochs\ epochs at a learning rate of
$\cpuHeadLr$, which gives \cpuTotalSteps\ steps and a budget of \cpuLrSum, with the other
settings unchanged: AdamW, no weight decay, linear decay, gradient clipping at one, dropout
$0.1$ and a batch of 64. These values are ours, they are listed here because they are not the
benchmark's, and they are identical for every method in the track, which is what the comparison
requires. Every method is run with five seeds, 42 to 46, and reported as a mean with a standard
deviation. The checkpoint selection rule, the losses, the decoders, the collapse tables and the
metrics are the same code as in the fine-tuning track.

This track also carries one ablation that the fine-tuning track cannot afford: CORAL with
zero-initialised cut-points. Here the head budget is \cpuLrSum, far more than the spread the
cut-points need, so the initialisation of Section~\ref{sec:meth-coral-init} should make little
difference; under the fine-tuning protocol the budget is \lrSumEncoder\ and it should make a
large one. The two tracks together therefore test the argument from both sides.

Because the encoder never adapts, the absolute numbers of this track are much lower than the
fine-tuning numbers, which is the expected outcome and not a defect. They are never compared
with the published tables.

\section{Implementation and Reproducibility}
\label{sec:meth-impl}

The fine-tuning track is a single self-contained script, run in three parts on three accounts
of a free notebook service so that the eleven runs could proceed in parallel; the part is only
a scheduling label and no hyper-parameter depends on it. Each run took between \minAll\ and \maxAll\ minutes on one NVIDIA Tesla T4, and the eleven runs together took \totalMinutes\
minutes of accelerator time, that is about \totalHours\ hours.

Every result file records the software and hardware it was produced on: Python
\pythonVersion, PyTorch \torchVersion, Transformers \tfVersion, Datasets \datasetsVersion, the
accelerator name, the number of visible accelerators, the realised training batch, the observed
number of optimiser steps per epoch, whether mixed precision was active, the corpus repository
revision (\datasetSha) and the encoder repository revision (\modelSha). Pinning the two
repository revisions matters because both are living repositories that can change after
publication.

Three guards run in every training run, and each aborts the run rather than produce a number
that cannot be trusted. The first refuses to start if more than one accelerator is visible or
if the realised training batch is not 64; an earlier version of these experiments silently
trained at a batch of 128 because a second accelerator was present, and this guard exists to
make that failure impossible. The second records the number of optimiser steps actually
performed in each epoch and checks it against the number implied by the corpus size, the
micro-batch and the accumulation. The third is the reloaded-checkpoint comparison of
Section~\ref{sec:meth-protocol}. The outcome of all three is reported for every method in
Chapter~\ref{chap:results}.

\section{Summary}

This chapter fixed everything except the head and the loss: the corpus and its official splits,
the D3Tok input read from the corpus itself, the AraBERTv2 encoder, and a training protocol
taken from the benchmark script down to the library defaults. Nine runs form the comparable
group, one two-phase run implements the metric-aligned objective, and two ablations test the
cut-point argument and the decision threshold, with a third ablation testing the checkpoint
selection rule. Every loss and every metric was verified against an independent implementation
before training, the three differences from the published script were declared, and the
cut-point initialisation of Equation~\ref{eq:coral-init} was introduced with an explicit
prediction about what the experiments should show. A second, frozen-backbone track provides
the CPU-friendly setting of the research brief and the variance estimates that a single-seed
protocol cannot give.

\textit{The next chapter reports what these eleven runs produced, checks them against the
published benchmark, and tests the predictions made here.}
